\subsection{机械臂控制}\label{sec:arm}

抓取这一步与底盘行进完全不同。底盘只需把两个轮速持续刷新给电机，机械臂却要把多个关节按
先后次序摆到一组确定的角度，再让末端的夹爪在合适的时机合拢或张开。这类动作天然由一块专用
的运动控制器在臂体本地闭环执行，它读取各关节的当前状态，按指令把关节驱到目标位姿并维持，
主控只负责下达"去哪个位姿"这样的高层意图。主控与这块控制器之间因此需要一条串行链路来往返
传递状态与指令，而 RK3588 主板对外并没有现成的串口引脚直连到臂上，机械臂一侧给出的是一个
USB 接口。中间的桥接由一颗 CP210x USB 转串口芯片承担，它在 USB 与异步串口之间做协议转换，
在系统里以一个串口设备的形态呈现，应用打开 \texttt{/dev/ttyUSB0} 即可像操作普通串口那样与
臂上的运动控制器通信，读取臂的状态并发出定位与抓放指令。

\repfig{robotic-arm.png}{多关节机械臂与夹爪，经 CP210x USB 转串口桥接挂到 RK3588 主板的 \texttt{/dev/ttyUSB0}}{fig:arm}

Starry 起初并没有这条 USB 串口通路。USB 主机栈枚举到 CP210x 设备后，并不会自动把它变成一个
可供用户态读写的串口节点，缺的是一层把 CP210x 的厂商控制传输翻译为串口语义、并把它登记为
\texttt{/dev/ttyUSB0} 的设备驱动。\textbf{本队为此在内核里补齐了 USB 串口 tty 支持，其中的 CP210x
后端对齐 Linux 同名驱动的行为，完成设备的初始化序列、波特率配置与控制传输处理，使枚举到的
CP210x 在 Starry 上呈现为标准的 \texttt{/dev/ttyUSB0}。}\textbf{这部分工作以合并请求
\href{https://github.com/rcore-os/tgoskits/pull/1378}{\#1378} 提交至上游并已合并，附带的主机端
单元测试覆盖了初始化与配置路径，因此这条通路不是临时打通的，而是按上游可维护的形态落地的。}

\paragraph{把多步位姿收敛为几个高层动作}
控制回路按帧产出决策，它没有余裕去逐关节地编排运动。把一次完整抓取展开来看，臂要先回到一个
张开的待命位姿，发现球后伸向目标、合拢夹爪、抬起，找到回收筐后再在筐口上方张开夹爪松手，
每一步都是若干关节角的有序变化。直接把这套序列摊在控制回路里，回路就会被关节级的细节淹没，
也会在等待臂体到位时阻塞,错过新到的帧。机械臂的后端抽象因此把这些多步序列归并为面向回路的
几个高层动词,分别是回到待命位姿、抓取、以及松手投放。待命动词把臂带回原点并张开夹爪；抓取
动词把伸向目标、合拢与抬起这一连串动作合为一体；松手动词在筐口上方张开夹爪完成投放。回路
只需在状态机给出抓或放的判定时下达对应的一个动词，余下的关节编排沿串口链路在臂上自行展开。

\paragraph{与控制回路的衔接}
这几个动词向控制回路呈现为非阻塞的调用，状态机每帧至多触发其一。控制器在收到一帧检测结果
并跑完状态机后，只在动作判定不为空时才按判定选择待命、抓取或松手中的一个下达，随即继续处理
下一帧，而不在调用点上等待臂体到位。这样抓放这类耗时较长的物理动作就不会顶在帧到指令的关键
路径上，感知与控制得以维持稳定的帧率，机械臂则沿 \texttt{/dev/ttyUSB0} 这条串口链路把收到的
高层动词落实为对实物臂的关节定位与夹爪开合，真实地完成对网球的抓取与在回收筐上方的投放。
